\documentclass[12pt]{article}
\usepackage[margin=0.7in]{geometry}

% \usepackage[margin=0.8in]{geometry}

% \usepackage{newpxtext}
% \usepackage{newpxmath}

\usepackage[tt=false]{libertine}

% \usepackage{newtxtext}
% \usepackage{newtxmath}

\usepackage[stretch=10]{microtype}

\linespread{1.05}

\usepackage{enumitem}
\setlist{nosep}

\usepackage[hidelinks, linktoc=all]{hyperref}

% Answers to the review questions. In the PDF this typesets the answer as an
% indented block right after the question; on the web the same environment is
% turned into a click-to-reveal dropdown (see web-pipeline/qa-details.lua).
% Pandoc does not implement \NewDocumentEnvironment, so it leaves the environment
% as a Div the filter can restructure, while LaTeX honours the definition here.
\NewDocumentEnvironment{answer}{}
  {\par\smallskip\noindent\ignorespaces}
  {\par\smallskip}

\title{Ethics}
\date{}
\author{}


\begin{document}

\section{Questions}
\subsection{Metaethics}
\begin{enumerate}
    \item Both non-cognitivists and error theorists agree that morality is an illusion because it does not exist. TRUE or FALSE.
    \begin{answer}\textbf{FALSE.} Non-cognitivism only denies that moral statements are truth-apt (they express emotion), and being a non-cognitivist does not make one a moral skeptic. Error theory says moral statements \emph{attempt} to state facts but are all false. It is Extreme Nihilism (an extension of error theory), not non-cognitivism, that calls morality an illusion --- so the two views do not share that claim.\end{answer}
    \item Cognitivists think that when moral judgments are true, they are always proven true in virtue of God’s will. TRUE or FALSE?
    \begin{answer}\textbf{FALSE.} Cognitivism only claims that moral statements are truth-apt; \emph{what} makes them true varies (e.g.\ natural facts under ethical naturalism). Grounding moral truth in God’s will is specifically Divine Command Theory / ethical supernaturalism, not a commitment of all cognitivists.\end{answer}
    \item According to Harman, we cannot use evidence in ethical observations to justify the truth of moral principles because moral principles are never true. TRUE or FALSE?
    \begin{answer}\textbf{FALSE.} Harman’s actual reason is that moral facts are explanatorily idle --- unnecessary to explain our observations (Ockham’s Razor), so we have no evidence for them. He does not claim moral principles are ``never true'' (that is error theory). The stated reasoning is wrong.\end{answer}
\end{enumerate}
\subsection{Moral Relativism}
\begin{enumerate}
    \item If you believe in (empirical) cultural relativism, then you must be a cognitivist about moral judgments. TRUE or FALSE?
    \begin{answer}\textbf{FALSE.} Descriptive (empirical) cultural relativism is just the observation that beliefs about right and wrong vary across cultures. It says nothing about whether moral judgments are truth-apt, so it carries no cognitivist commitment.\end{answer}
    \item Utilitarians are not moral relativists because…
    \begin{enumerate}
        \item The utility principle applies to everyone regardless of their culture.
        \item Benthamic utilitarians are often interpreted as ethical naturalists.
        \item There can be only one universally right action for every moral predicament.
        \item All of the above.
        \item a and c only.
    \end{enumerate}
    \begin{answer}\textbf{(d) All of the above.} The utility principle is universal and impartial (a); Benthamic utilitarians are ethical naturalists, who hold that objective moral facts exist --- the opposite of relativism (b); and it posits a single universal standard that fixes the right action (c). All three make utilitarianism universalist rather than relativist.\end{answer}
\end{enumerate}

\subsection{Utilitarianism}
\begin{enumerate}
    \item The principle of utility requires that we always perform the action that maximizes well-being, therefore, it entails that we ought to always choose the action that benefits the greatest number of people. TRUE or FALSE?
    \begin{answer}\textbf{FALSE.} The utility principle maximizes \emph{total net utility}, not head-count. A huge benefit to two people can outweigh a tiny benefit to a hundred, so ``the greatest number of people'' does not follow.\end{answer}
    \item The following is an example of an agent-relative reason
    \begin{enumerate}
        \item Every parent ought to save their own children because parents have special obligation to their children.
        \item Everyone has a duty to save a drowning child when it is a very easy rescue.
        \item Everyone has a duty to save people when they themselves are feeling happy.
        \item Every medical doctor ought not to kill an innocent patient even if he can save more people because medical doctors have made promises to do no harm.
        \item Options a, b and d
        \item Options a and d
    \end{enumerate}
    \begin{answer}\textbf{(f) Options a and d.} Both are universal rules applied to a specific relationship or role (a parent to their own child; a doctor bound by their promise), which is what makes a reason agent-relative. (b) is an agent-neutral duty that holds for everyone; (c) is a subjective feeling, not an agent-relative moral rule.\end{answer}
    \item Utilitarianism is a defective theory because we cannot always predict the consequences of our actions, so it cannot help us deliberate and make decisions about what to do. (Note: True/False based on how utilitarians answer the Deliberation objection)
    \begin{answer}\textbf{FALSE.} Utilitarians answer the Deliberation objection by separating the \emph{Standard of Rightness} (what makes an act correct) from the \emph{Decision Procedure} (how we actually choose --- proven rules of thumb, the basketball analogy). So the objection does not show the theory is defective.\end{answer}
\end{enumerate}

\subsection{Deontology}
\begin{enumerate}
    \item Kantian deontologists argue that we have a perfect duty to help others. TRUE or FALSE?
    \begin{answer}\textbf{FALSE.} Helping others is an \emph{imperfect} duty (flexible in when, how, and toward whom it is fulfilled), not a perfect one.\end{answer}
    \item Under Kant's Principle of Humanity, it is always immoral to use another human being as a means to an end. TRUE or FALSE?
    \begin{answer}\textbf{FALSE.} It is immoral only to use someone as a \emph{mere} means. Using another as a means is permissible when they share the purpose (the barista example). The word ``mere'' is the whole point.\end{answer}
    \item According to Kant, imperfect duties are less important than perfect duties because imperfect duties are merely optional suggestions, not actual moral obligations. TRUE or FALSE?
    \begin{answer}\textbf{FALSE.} Imperfect duties are still strict moral requirements: you \emph{must} adopt the goal (e.g.\ helping others). Only the when, how, and toward whom is left to your discretion --- they are not optional suggestions.\end{answer}
    \item According to the principle of humanity, we ought to treat humans never as mere means, but always as ends in themselves. This means that it is morally wrong to be the CEO of any company because I will be making people work for me. TRUE or FALSE?
    \begin{answer}\textbf{FALSE.} Employees share the purpose of the arrangement (like the barista getting paid for the transaction), so they are not being used as a \emph{mere} means. Employing people is not a violation of the Principle of Humanity.\end{answer}
    \item Based on Ellis’ paper, deontologists can appeal to legal prohibitions to explain the idea of “degree wrongness” but this can be problematic because…
    \begin{enumerate}
        \item The idea of intrinsic wrongness is essential to the deontologists.
        \item Even if we accept the very idea of degree of wrongness, we still cannot justify the cut-off point without being arbitrary.
        \item It is not the case killing is always the worst wrongdoing even though we tend to take the duty to not kill much more seriously than other duties.
        \item All of the above.
        \item b and c only.
    \end{enumerate}
    \begin{answer}\textbf{(d) All of the above.} Ellis’ Legal Model Trap: deontology rests on \emph{intrinsic}, reason-based wrongness (a), but legal categories are a man-made convention whose cut-offs are arbitrary (b) and whose severity does not track intrinsic moral wrongness (c). Borrowing the legal framework therefore reduces universal morality to human convention.\end{answer}
\end{enumerate}

\end{document}
